\Chapter{4}

\Verse{1}
{
	\hJ{והאדם ידע את חוה אשתו ותהר ותלד את קין ותאמר קניתי איש את יהוה}
}
{
	\eJ{
		And the man knew Eve his wife,\\
		and she conceived and bore Cain,\\
		and said, “I have acquired a man with YHWH.”\\
	}
}
{
}

\Verse{2}
{
	\hJ{ותסף ללדת את אחיו את הבל ויהי הבל רעה צאן וקין היה עבד אדמה}
}
{
	\eJ{
		And she continued to bear his brother, Abel;\\
		and Abel became a keeper of sheep,\\
		and Cain was a worker of the ground.\\
	}
}
{
}

\Verse{3}
{
	\hJ{ויהי מקץ ימים ויבא קין מפרי האדמה מנחה ליהוה}
}
{
	\eJ{
		And it came to pass, after some days,\\
		that Cain brought from the fruit of the ground\\
		an offering to YHWH.\\
	}
}
{
}

\Verse{4}
{
	\hJ{והבל הביא גם הוא מבכרת צאנו ומחלבהן וישע יהוה אל הבל ואל מנחתו}
}
{
	\eJ{
		And Abel, he also brought from the firstborn of his flock\\
		and from their fat portions;\\
		and YHWH regarded Abel and his offering.\\
	}
}
{
}

\Verse{5}
{
	\hJ{ואל קין ואל מנחתו לא שעה ויחר לקין מאד ויפלו פניו}
}
{
	\eJ{
		But to Cain and to his offering He did not regard;\\
		and Cain burned greatly with anger,\\
		and his face fell.\\
	}
}
{
}

\Verse{6}
{
	\hJ{ויאמר יהוה אל קין למה חרה לך ולמה נפלו פניך}
}
{
	\eJ{
		And YHWH said to Cain,\\
		“Why are you angry,\\
		and why has your face fallen?”\\
	}
}
{
}

\Verse{7}
{
	\hJ{הלוא אם תיטיב שאת ואם לא תיטיב לפתח חטאת רבץ ואליך תשוקתו ואתה תמשל בו}
}
{
	\eJ{
		“Is it not so, if you do well, you will be lifted up?\\
		And if you do not do well,\\
		sin is crouching at the door;\\
		and its desire is toward you,\\
		but you must rule over it.”\\
	}
}
{
}

\Verse{8}
{
	\hJ{ויאמר קין אל הבל אחיו ויהי בהיותם בשדה ויקם קין אל הבל אחיו ויהרגהו}
}
{
	\eJ{
		And Cain said to Abel his brother;\\
		and it came to pass, when they were in the field,\\
		that Cain rose up against Abel his brother\\
		and killed him.\\
	}
}
{
}

\Verse{9}
{
	\hJ{ויאמר יהוה אל קין אי הבל אחיך ויאמר לא ידעתי השמר אחי אנכי}
}
{
	\eJ{
		And YHWH said to Cain,\\
		“Where is Abel your brother?”\\
		And he said, “I do not know;\\
		am I my brother’s keeper?”\\
	}
}
{
}

\Verse{10}
{
	\hJ{ויאמר מה עשית קול דמי אחיך צעקים אלי מן האדמה}
}
{
	\eJ{
		And He said,\\
		“What have you done?\\
		The voice of your brother’s blood\\
		cries out to me from the ground.\\
	}
}
{
}

\Verse{11}
{
	\hJ{ועתה ארור אתה מן האדמה אשר פצתה את פיה לקחת את דמי אחיך מידך}
}
{
	\eJ{
		“And now you are cursed from the ground,\\
		which has opened its mouth\\
		to receive your brother’s blood from your hand.\\
	}
}
{
}

\Verse{12}
{
	\hJ{כי תעבד את האדמה לא תסף תת כחה לך נע ונד תהיה בארץ}
}
{
	\eJ{
		“When you work the ground,\\
		it shall no longer yield its strength to you;\\
		a fugitive and a wanderer\\
		you shall be on the earth.”\\
	}
}
{
}

\Verse{13}
{
	\hJ{ויאמר קין אל יהוה גדול עוני מנשא}
}
{
	\eJ{
		And Cain said to YHWH,\\
		“My punishment is greater than I can bear.\\
	}
}
{
}

\Verse{14}
{
	\hJ{הן גרשת אתי היום מעל פני האדמה ומפניך אסתר והייתי נע ונד בארץ והיה כל מצאי יהרגני}
}
{
	\eJ{
		“Behold, you have driven me out today\\
		from the face of the ground,\\
		and from your presence I shall be hidden;\\
		and I shall be a fugitive and a wanderer on the earth,\\
		and whoever finds me will kill me.”\\
	}
}
{
}

\Verse{15}
{
	\hJ{ויאמר לו יהוה לכן כל הרג קין שבעתים יקם וישם יהוה לקין אות לבלתי הכות אתו כל מצאו}
}
{
	\eJ{
		And YHWH said to him,\\
		“Therefore, whoever kills Cain,\\
		vengeance shall be taken on him sevenfold.”\\
		And YHWH set a mark for Cain,\\
		so that none who found him would strike him.\\
	}
}
{
}

\Verse{16}
{
	\hJ{ויצא קין מלפני יהוה וישב בארץ נוד קדמת עדן}
}
{
	\eJ{
		And Cain went out from the presence of YHWH
		and dwelt in the land of Nod, east of Eden.
	}
}
{

\begin{quote}
\textbf{נוד} (\textit{nód}) \emph{m.}

Meaning \emph{wanderings}, \emph{vagabondage}, or \emph{journey}. Occurs without plural forms; in the construct state (\textit{נוֹד־}), conveys a sense of \emph{migration} or \emph{migratory movement}.
\end{quote}

\aA{So he goes to Iran.}

\aB{Iran?}

\aA{Yeah it's east of eden.}

}


\Verse{17}
{
	\hJ{וידע קין את אשתו ותהר ותלד את חנוך ויהי בנה עיר ויקרא שם העיר כשם בנו חנוך}
}
{
	\eJ{
		And Cain knew his wife,\\
		and she conceived and bore Enoch;\\
		and he became a builder of a city,\\
		and called the name of the city
		after the name of his son, Enoch.\\
	}
}
{
}

\Verse{18}
{
	\hJ{ויולד לחנוך את עירד ועירד ילד את מחויאל ומחויאל ילד את מתושאל ומתושאל ילד את למך}
}
{
	\eJ{
		And to Enoch was born Irad;\\
		and Irad begot Mehujael;\\
		and Mehujael begot Methushael;\\
		and Methushael begot Lamech.\\
	}
}
{

\aB{


}

}

\Verse{19}
{
	\hJ{ויקח לו למך שתי נשים שם האחת עדה ושם השנית צלה}
}
{
	\eJ{
		And Lamech took for himself two wives:\\
		the name of the one was Adah,\\
		and the name of the other, Zillah.\\
	}
}
{
}

\Verse{20}
{
	\hJ{ותלד עדה את יבל הוא היה אבי ישב אהל ומקנה}
}
{
	\eJ{
		And Adah bore Jabal;\\
		he was the father of those who dwell in tents\\
		and have livestock.\\
	}
}
{
}

\Verse{21}
{
	\hJ{ושם אחיו יובל הוא היה אבי כל תפש כנור ועוגב}
}
{
	\eJ{
		And the name of his brother was Jubal;\\
		he was the father of all those who play the lyre\\
		and the pipe.\\
	}
}
{
}

\Verse{22}
{
	\hJ{וצלה גם הוא ילדה את תובל קין לטש כל חרש נחשת וברזל ואחות תובל קין נעמה}
}
{
	\eJ{
		And Zillah, she also bore Tubal-Cain,\\
		the forger of every cutting instrument\\
		of bronze and iron;\\
		and the sister of Tubal-Cain was Naamah.\\
	}
}
{

}

\Verse{23}
{
	\hJ{ויאמר למך לנשיו עדה וצלה שמען קולי נשי למך האזנה אמרתי כי איש הרגתי לפצעי וילד לחברתי}
}
{
	\eJ{
		And Lamech said to his wives,\\
		“Adah and Zillah, hear my voice;\\
		wives of Lamech, listen to what I say:\\
		for I have killed a man for wounding me,\\
		even a young man for striking me.”\\
	}
}
{

}


\Verse{24}
{
	\hJ{כי שבעתים יקם קין ולמך שבעים ושבעה}
}
{
	\eJ{
		If Cain is avenged sevenfold,\\
		then Lamech seventy-sevenfold.\\
	}
}
{

%%%%%%%%%%%%%%%%%
%%% ATTENTION %%%
%%%%%%%%%%%%%%%%%
%
% The Paleo Hebrew strings are REVERSED here.
%
% That's on purpose, for now, but this needs
% to be fixed with a proper \Paleo command.


\aA{

	\begin{center}
	\begin{tabular}{lr}

	Irad & {\Paleo 𐤃𐤓𐤉𐤏} \\[1.5em]

	his city & {\Paleo 𐤅𐤓𐤉𐤏} \\[1.5em]

	\end{tabular}	
	\end{center}

	\begin{center}
	\begin{tabular}{lr}

	Jabal & {\Paleo 𐤋𐤁𐤉} \\[1.5em]

	Jubal & {\Paleo 𐤋𐤁𐤅𐤉} \\[1.5em]

	Tubal & {\Paleo 𐤋𐤁𐤅𐤕} \\[1.5em]

	Cain & {\Paleo 𐤍𐤉𐤒} \\[1.5em]
	
	\end{tabular}	
	\end{center}

	A mark for Cain.

	\begin{center}
	\begin{tabular}{lr}

	Mehujael & {\Paleo 𐤋𐤀𐤉𐤅𐤇𐤌} \\[1.5em]

	Methushael & {\Paleo 𐤋𐤀𐤔𐤅𐤕𐤌} \\[1.5em]

	\end{tabular}	
	\end{center}

	\begin{center}
	\begin{tabular}{lr}
	
	Adah & {\Paleo 𐤄𐤃𐤏} \\[1.5em]

	Zillah & {\Paleo 𐤄𐤋𐤑} \\[1.5em]

	\end{tabular}	
	\end{center}

}

	\aB{I don't see any pattern there.}

	\aA{It's J.}

	\aB{Meaning what?}

	\aA{J is \emph{never} not up to something.}

\aA{

	\begin{quote}
	\textbf{Adah} (\textit{Hebrew}: \hAJ{עדה}) \emph{n.}\\
	A group or meeting (particularly for purposes of worship or judicial matters);\\
	congregation; assembly; prayer meeting; a gathering; a company assembled together.
	See also \hAJ{עִדּוּי} conception, pregnancy. From Aram., verbal n. of \hAJ{עַדִּי} (she conceived, became pregnant).


	\medskip
	
	\textbf{Zillah} (\textit{Hebrew}: \hAJ{צלה}) \emph{n.}\\
	Cognate with Aramaic \textit{\hAJ{צלא}} (“to pray,” “to bend down in prayer”) and Arabic \textit{ṣalā} (“to roast,” “to burn”).\\
	Also associated with \textit{\hAJ{צלע}} (\textit{tselá}), meaning \emph{rib} (of the body), \emph{rib} (of meat),
	or \emph{rib} (in architectural contexts).
	\end{quote}

	% Adah means a group or meeting (particularly for purposes of worship or judicial matters); congregation; assembly; prayer meeting

	% Zillah is cognate with Aramaic צלא which means to pray, to bend down (to pray); and Arabic صَلَى (ṣalā, “to roast, fry”).
	% Zillah also is cognate to a word meaning "rib"
	%
	% Def: צֵלָע • (tselá) f (plural indefinite צְלָעוֹת or צְלָעִים, singular construct צֶלַע־ or צֵלַע־, plural construct צַלְעוֹת־ or צַלְעֵי־) [pattern: קֵטָל]
	% rib (of the body)
	% rib (of meat)
	% rib (of a door or leaf, a bent or curved architectural supporting element)

	% https://en.wiktionary.org/wiki/%D7%A2%D7%93%D7%94
	% https://en.wiktionary.org/wiki/%D7%A6%D7%9C%D7%90#Aramaic
	% https://en.wiktionary.org/wiki/%D7%A6%D7%9C%D7%A2

	\aA{Sometimes the signals aren't clear.}


	\begin{center}
	\begin{tabular}{lcr}
	Jabal & \hAJ{יבל} & yā·ḇāl \\
	he & \hAJ{הוא} & hū \\
	was & \hAJ{היה} & hā·yāh \\
	the father & \hAJ{אבי} & ’ă·ḇî \\
	of those who dwell & \hAJ{ישב} & yō·šēḇ \\
	in tents & \hAJ{אהל} & ’ō·hel \\
	and [raise] livestock & \hAJ{ומקנה} & ū·miq·neh \\
	\end{tabular}
	\end{center}

	\aA{You've got to check everything with J.}

	\aB{Sounds hard.}

	\aA{It's not.}

	\begin{center}
	\begin{tabular}{lr}

	seven & {\Paleo 𐤏𐤁𐤔} \\[1.5em]

	seventy & {\Paleo 𐤌𐤉𐤏𐤁𐤔} \\[1.5em]

	sevenfold & {\Paleo 𐤌𐤉𐤕𐤏𐤁𐤔} \\[1.5em]

	seventy & {\Paleo 𐤌𐤉𐤏𐤁𐤔} \\[0.25em]
	-seven  & {\Paleo 𐤄𐤏𐤁𐤔}  \\[0.25em]

	\end{tabular}	
	\end{center}

}


}


\Verse{25}
{
	\hR{וידע אדם עוד את אשתו ותלד בן ותקרא את שמו שת כי שת לי אלהים זרע אחר תחת הבל כי הרגו קין}
}
{
	\eR{
		And Adam knew his wife again,\\
		and she bore a son\\
		and called his name Seth:\\
		“For God has appointed me another seed\\
		in place of Abel, for Cain killed him.”\\
	}
}
{
}

\Verse{26}
{
	\hR{ולשת גם הוא ילד בן ויקרא את שמו אנוש}
	\hJ{אז הוחל לקרא בשם יהוה}
}
{
	\eR{
		And to Seth, to him also a son was born,\\
		and he called his name Enosh;
	}
	\eJ{
		Then it was begun to call on the name YHWH.\\
	}
}
{

}
